\chapter*{Mở đầu}
\addcontentsline{toc}{chapter}{Mở đầu}
\setcounter{page}{1}
\pagenumbering{arabic}
Trong kỷ nguyên công nghiệp 4.0, sự bùng nổ của thị trường thiết bị di động đã đặt ra những yêu cầu ngày càng khắt khe đối với quy trình kiểm soát chất lượng (Quality Control - QC) trong các dây chuyền sản xuất. Màn hình điện thoại, với vai trò là giao diện tương tác chính giữa người và máy, là thành phần đòi hỏi độ hoàn thiện cao nhất. Bất kỳ một khiếm khuyết nhỏ nào về hiển thị, từ các điểm chết (dead pixels) đến các đường kẻ sọc hay vết loang màu, đều ảnh hưởng trực tiếp đến trải nghiệm người dùng và uy tín của nhà sản xuất. Do đó, việc phát triển các hệ thống kiểm tra quang học tự động (AOI) với độ chính xác cao và tốc độ xử lý thời gian thực đang trở thành một nhu cầu cấp thiết nhằm thay thế hoàn toàn sự phụ thuộc vào mắt thường của con người.

Tuy nhiên, các giải pháp thị giác máy tính truyền thống và các mô hình học sâu có giám sát (Supervised Deep Learning) hiện nay đang đối mặt với những rào cản lớn khi triển khai thực tế. Các phương pháp này đòi hỏi một lượng dữ liệu khổng lồ được gán nhãn thủ công tỉ mỉ cho từng loại lỗi, điều này tốn kém cả về thời gian và chi phí nhân lực. Hơn nữa, trong môi trường sản xuất năng động, các dạng lỗi mới thường xuyên xuất hiện một cách bất ngờ và không thể dự báo trước. Các mô hình mạng nơ-ron tĩnh (static models) thường thất bại trong việc nhận diện các mẫu lỗi lạ này, hoặc nếu được huấn luyện lại trên dữ liệu mới, chúng lại mắc phải hiện tượng "quên thảm họa" (Catastrophic Forgetting) – tức là mất đi khả năng nhận diện các lỗi cũ đã học trước đó.

Để giải quyết những hạn chế nêu trên, hướng tiếp cận sử dụng Học máy không giám sát (Unsupervised Learning) kết hợp với Học liên tục (Continual Learning) đang nổi lên như một giải pháp đầy tiềm năng. Thay vì dựa vào nhãn dữ liệu có sẵn, hệ thống có thể tự học cách biểu diễn các đặc trưng của dữ liệu bình thường và phát hiện các bất thường dựa trên sự sai lệch phân bố. Đặc biệt, việc ứng dụng các thuật toán có khả năng thích nghi cấu trúc mạng (topology adaptation) cho phép mô hình linh hoạt mở rộng "vùng tri thức" của mình theo thời gian, tiếp nhận các loại lỗi mới mà không làm xáo trộn những tri thức nền tảng đã được thiết lập.

Luận văn này đề xuất một kiến trúc hệ thống lai ghép (hybrid architecture) tiên tiến nhằm kiểm tra lỗi màn hình tự động. Giải pháp kết hợp sức mạnh trích xuất đặc trưng của mạng Autoencoder nén dữ liệu với khả năng phân cụm động của mạng Growing Neural Gas (GNG). Điểm mới của nghiên cứu nằm ở việc tích hợp kênh thông tin Entropy cục bộ (Local Entropy) vào quá trình tiền xử lý, giúp tăng cường khả năng phát hiện các lỗi có độ tương phản thấp. Đồng thời, cơ chế sinh trưởng nơ-ron và cạnh tranh của GNG được khai thác triệt để nhằm giải quyết bài toán học liên tục, đảm bảo hệ thống có thể nhận diện chính xác các lỗi hình học phức tạp như hình tròn, hình tam giác, hay hình chữ nhật qua từng giai đoạn sản xuất mà không cần huấn luyện lại từ đầu.

Mục tiêu chính của luận văn là xây dựng và kiểm chứng một quy trình kiểm tra lỗi hoàn chỉnh, từ khâu sinh dữ liệu giả lập, trích chọn đặc trưng đến phân loại thích nghi. Các đóng góp của đề tài không chỉ dừng lại ở việc cải thiện độ chính xác phát hiện lỗi mà còn cung cấp một minh chứng thực nghiệm rõ ràng cho khả năng khắc phục hiện tượng quên thảm họa trong các hệ thống trí tuệ nhân tạo công nghiệp. Kết quả nghiên cứu hứa hẹn mở ra hướng đi mới cho việc phát triển các hệ thống giám sát chất lượng thông minh, linh hoạt và tiết kiệm chi phí cho các nhà máy thông minh trong tương lai.

Về bố cục, nội dung của luận văn bao gồm 4 Chương như sau: Chương \ref{chap:c1} đưa ra bài toán phát hiện các lỗi hiển thị bất thường trên màn hình điện thoại. Chương \ref{chap:c2} nền tảng kiến thức cơ bản, các nghiên cứu liên quan đến học không giám sát thích ứng và học động. Chương \ref{chap:c3} đưa ra phương pháp phát hiện bất thường bằng học không giám sát thích ứng để giải quyết bài toán. Chương \ref{chap:c4} kết quả thực nghiệm bằng ứng dụng thực tế sinh dữ liệu và học liên tục với mạng GNG.Cuối cùng là Kết luận tổng kết lại kết quả mà luận văn đạt được và đề xuất hướng phát triển. 
\cleardoublepage